\section{Repository audit and the missing geometric layer}

\subsection{Pinned repository state}

The mathematical audit is pinned to commit
\code{bce5113f44d613f5625c274128ee761f35486fc5}.  During final verification, \code{main} had advanced to
\code{a2f72b0cdf57f8a3fefd34b272c3d471385e819c}, but a direct contents-API comparison showed that the three entries under
\code{Logic/PrimitiveRecursiveArithmetic} were unchanged: \code{ConciseIPRA} had Git object
\code{622927389849d1e4c3b5f813b8481b7674974acc}, \code{README.md} had object
\code{1340d58e7f9c97cc8593341f5e80452633a2b77d}, and \code{docs} had object
\code{996363773c5e9f09fa826e0cc943f55af059d8dd} at both commits.  Thus unrelated repository activity did not alter the
audited iPRA subtree.

At the pinned snapshot the root iPRA README listed sixty-six research packets.  It also stated an
important caveat: the packets were filed as delivered and were not independently rerun, compiled,
or verified by the indexing commit.  Consequently this report treats packet claims as research
inputs, not as automatically certified facts.

The checked foundation is \code{ConciseIPRA}.  The relevant Lean layers are:

\begin{center}
\begin{tabularx}{0.94\textwidth}{@{}lX@{}}
\toprule
Module/interface & Role \\
\midrule
\code{Syntax.lean} & First-order terms and formulas, substitution, and quantifier-free formulas. \\
\code{Deduction.lean} & Intuitionistic natural deduction with equality; derivations \code{Prf}. \\
\code{PR.lean} & The inductive family \code{PR}, terms \code{PTm}, axioms \code{PRAAx}, and the quantifier-free induction interface. \\
\code{NativeBiInterpretation.lean} & The bridge \code{pra_iff_native}. \\
\bottomrule
\end{tabularx}
\end{center}

This package therefore targets the usual iPRA calculus first.  Native export is the last step, not a
substitute for an embedded derivation.

\subsection{Prior frameworks that are treated as baseline}

The repository already contains a substantial general methodology.  The following table records the
parts most relevant here.

\begin{longtable}{@{}p{0.25\textwidth}p{0.31\textwidth}p{0.36\textwidth}@{}}
\toprule
Packet & Existing contribution & Boundary retained here \\
\midrule
\endfirsthead
\toprule
Packet & Existing contribution & Boundary retained here \\
\midrule
\endhead
\code{ipra_framework} & Finite contracts, bounded Booleanization, bounded search, quotient-free finite presentations, and a compiler outline. & A contract does not by itself provide geometric overlap objects, localization equality, or a sheaf gluing formula. \\
\code{ipra_embedding_framework} & Demand-driven finite observations, certificate-relative interpretation, finite windows, simple Cauchy branches, and finite algebra interfaces. & Exact global branch data and general height bounds remain separate producer problems. \\
\code{ipra_finite_observation} & Precise diagnosis of illegal universal evaluation, Cauchy-induction mistakes, computation/proof horizons, and exact-output demands. & All algorithms below are fixed PR constructions; no universal interpreter is introduced. \\
\code{ipra_coherent_methods} & Coherent finite presentations and several small genuine \code{Prf PRAAx} derivations. & The weighted and object-specific applications are not automatically embedded. \\
\code{ipra_local_global} & Variable-length recursion and numerical local-to-global rules. & Zariski localization and module descent require additional algebraic witness calculus. \\
\code{ipra_finite_use_framework} & Relational finite-use composition, backward viability, branch-invariant outputs, robust covers, and root clusters. & The present package uses that philosophy but adds exact localization, unit-ideal covers, and sheaf equalizers. \\
\code{ipra_stable_observables} & Exact or stable observables, matrix certificates, and impossibility of complete effective bound certificates in general. & A certificate verifier is not a certificate producer; this distinction remains explicit. \\
\code{diophantine_geometry_ipra} & Rational projective point codes, heights, conics, local obstructions, explicit descent, and conditional finite generation. & It treats points and heights rather than a general scheme/atlas/sheaf layer. \\
\bottomrule
\end{longtable}

The most recent finite-use packet is especially close in spirit.  It allows a producer to return a
finite relation of locally valid answers and lets a bounded consumer choose a viable transcript.
That is useful for chart selection and root clusters.  It does not, however, supply the algebraic
identity that makes a principal-open cover jointly monic or the formula that glues compatible module
fractions.  Those are the central additions here.

\subsection{Difficulties recorded by theorem-specific packets}

The theorem packets identify several logically distinct obstacles.

\begin{enumerate}[label=\arabic*.]
\item \textbf{No universal PR evaluator.}  An inductive Lean datatype of PR programs does not give
      an iPRA function evaluating arbitrary program codes.  Every evaluator in this package is fixed
      syntactically, or a finite trace is supplied and checked.
\item \textbf{No induction on infinitary predicates.}  Cauchy convergence, sheafhood over all opens,
      or ideal equality by unbounded search cannot be the induction formula.  Finite prefix folds and
      explicit witnesses replace them.
\item \textbf{Exact choice can be discontinuous.}  Singular Puiseux branches and arbitrary Cauchy
      coefficient zero tests remain outside a uniform exact interface.  Relational chart data or
      supplied separation witnesses are required.
\item \textbf{Cardinality does not locate points.}  The Faltings and Siegel packets stress that a
      bound on the number of rational or integral points is not a height cutoff and does not terminate
      enumeration.
\item \textbf{Descent requires actual representatives.}  The Mordell--Weil packet can consume a
      finite list covering a quotient but does not manufacture rational representatives from abstract
      Selmer data.
\item \textbf{Native Lean is not object iPRA.}  Several packets check host-level finite arithmetic
      but do not construct the corresponding \code{Prf PRAAx} value.
\end{enumerate}

The geometric framework below resolves none of the genuinely missing height or representative
producer theorems.  It does resolve a different structural difficulty: how to represent quotient
rings, localizations, finite covers, gluing, and projective cohomology without silently using
classical quotient types, prime spectra, arbitrary choice, or induction beyond iPRA.

\subsection{Why the Wikipedia checklist is useful but insufficient}

The user-specified overview divides algebraic geometry into polynomial zero sets, affine and
projective varieties, computational methods, and the modern scheme/sheaf viewpoint.  Each item
requires a different iPRA representation:

\begin{itemize}
\item polynomial zero sets use finite polynomial evaluation and bounded point searches;
\item affine geometry uses finitely presented coordinate rings;
\item projective geometry uses homogeneous presentations and finite standard charts;
\item computational geometry uses certificate-checked reductions rather than an assumed terminating
      universal algorithm;
\item schemes and sheaves use finite affine atlases and descent equalizers rather than an internally
      completed topology of prime ideals.
\end{itemize}

The rest of the report supplies one common architecture for these representations.
